\documentclass{amsart}
% 箇条書き
\usepackage{enumitem}

\title{Scalarization methods with a characterizing function 
of an ordering cone in set optimization}

\author{Ryota Iwamoto*}
\address[Ryota Iwamoto]
{Graduate School of Science and Technology, Niigata University,
8050 Ikarashi-2nocho, Niigata-City 950-2181, Japan}
\email{iwamoto@m.sc.niigata-u.ac.jp}
\thanks{*Presenting author}

\author{Tamaki Tanaka}
\address[Tamaki Tanaka]
{Department of Mathematics, Faculty of Science, Niigata University,
Niigata 950–2181, Japan.}
\email{tamaki@math.sc.niigata-u.ac.jp}

\keywords{set optimization; Fan-Takahashi minimax inequality; Hiriart-
Urruty oriented distance}

\begin{document}

\maketitle

\section{Introduction}
In convex analysis and optimization theory,
Fan-Takahashi minimax inequality plays a key role to solve equilibrium problems and
other problems. Let $X$ be a nonempty compact convex subset of a Hausdorff topological vector space and $f \colon X \times X \to \mathbb{R}$.
Fan-Takahashi minimax inequality is: if $f$ satisfies
the following conditions:
\begin{enumerate}[label=(\alph*)]
  \item for each fixed $y \in X$, $f(\cdot,y)$ is lower semicontinuous,
  \item for each fixed $x \in X$, $f(x,\cdot)$ is quasi concave,
  \item $f(x,x) \leq 0$ for all $x \in X$,
\end{enumerate}
then there exists $\bar{x} \in X$ such that $f(\bar{x},y) \leq 0$ for all $y \in X$. 
In addition, a number of researchers, including Georgiev and Tanaka, 
study Fan-Takahashi minimax inequalities for set-valued maps. In this talk,
we explain set-valued Fan-Takahashi minimax inequalities generalized by scalarization functions,
and consider relations between scalarization functions for set-valued Fan-Takahashi minimax inequalities 
and characterizing functions of the ordering cone $C$ \cite{EJ2021} defined by
\begin{equation}
  f(y) \leq 0 \Leftrightarrow y \in -C \quad \text{for all}\: y \in Y, \notag
\end{equation}
where $C$ is a convex cone in a real topological vector space $Y$ 
and $f: Y \to \mathbb{R} \cup \{\pm \infty\}$.

\begin{thebibliography}{88}
  \bibitem{EJ2021}
	G.\ Eichfelder and J.\ Jahn,
	{\em Optimality conditions in discrete-continuous nonlinear optimization}, 
	Minimax Theory Appl.
	{\bf 6} (2021), 127--144.
\end{thebibliography}

\end{document}